\begin{tikzpicture}
\node[minimum width=15cm,minimum height=9cm] (current page) {};
\begin{scope}[overlay]

\coordinate (box place) at ([yshift=0.25cm]current page.center);
\tikzset{
    box/.style={at=(box place), anchor=south,align=left,fill=white,draw=red,ultra thick}
}
\begin{visibleenv}<2>
\node[box] {space out accesses by 4096 \\
ensure separate cache sets and \\
avoid triggering prefetcher
};
\end{visibleenv}
\begin{visibleenv}<3>
\node[box] {
    repeat access if zero \\
    apparently value of zero speculatively read \\
    when real value not yet available
};
\end{visibleenv}
\begin{visibleenv}<4>
\node[box] {
    access cache to allow measurement later \\
    in paper with FLUSH+RELOAD instead \\\
    of PRIME+PROBE technique
};
\end{visibleenv}
\begin{visibleenv}<5>
\node[box] {
    segfault actually happens eventually \\
    option 1: okay, just start a new process every time \\
    (don't even need to hide segfault) \\
    option 2: way of suppressing segfault \\
    (paper used transactional memory support, \\
    not mispredicted branch)
};
\end{visibleenv}
\end{scope}
\end{tikzpicture}
